\documentclass[xcolor=dvipsnames,aspectratio=169]{beamer}

\usepackage{fontspec}
\defaultfontfeatures{Ligatures=TeX}
%\setsansfont{Liberation Sans}
\usepackage{polyglossia}
\setdefaultlanguage{german}

\usetheme{Berlin}
\usecolortheme[named=LimeGreen]{structure}
\usepackage{beamerthemesplit} % kam neu dazu
	
\usepackage{amsmath}
\usepackage{graphicx}
\graphicspath{{pictures/}}
\usepackage{amssymb}
\usepackage{amsfonts}
\usepackage{caption}
\usepackage{multimedia}
\usepackage{tikz}
\usepackage{listings}
\usepackage{acronym}
\usepackage{uhrzeit}
\usepackage{lmodern}
\usepackage{multicol}


\definecolor{pblue}{rgb}{0.13,0.13,1}
\definecolor{pgreen}{rgb}{0,0.5,0}
\definecolor{pred}{rgb}{0.9,0,0}
\definecolor{pgrey}{rgb}{0.46,0.45,0.48}

\newcounter{countitems}
\newcounter{nextitemizecount}
\newcommand{\setupcountitems}{%
  \stepcounter{nextitemizecount}%
  \setcounter{countitems}{0}%
  \preto\item{\stepcounter{countitems}}%
}
\makeatletter
\newcommand{\computecountitems}{%
  \edef\@currentlabel{\number\c@countitems}%
  \label{countitems@\number\numexpr\value{nextitemizecount}-1\relax}%
}
\newcommand{\nextitemizecount}{%
  \getrefnumber{countitems@\number\c@nextitemizecount}%
}
\newcommand{\previtemizecount}{%
  \getrefnumber{countitems@\number\numexpr\value{nextitemizecount}-1\relax}%
}

\makeatother    
\newenvironment{AutoMultiColItemize}{%
\ifnumcomp{\nextitemizecount}{>}{3}{\begin{multicols}{2}}{}%
\setupcountitems\begin{itemize}}%
{\end{itemize}%
\unskip\computecountitems\ifnumcomp{\previtemizecount}{>}{3}{\end{multicols}}{}}

\lstset{
    escapeinside={(*}{*)}
}

\lstdefinestyle{Java}{
  showspaces=false,
  showtabs=false,
  tabsize=2,
  breaklines=true,
  showstringspaces=false,
  breakatwhitespace=true,
  commentstyle=\color{pgreen},
  keywordstyle=\color{pblue},
  stringstyle=\color{pred},
  basicstyle=\footnotesize\ttfamily,
  numbers=left,
  numberstyle=\tiny\color{gray}\ttfamily,
  numbersep=7pt,
  %moredelim=[il][\textcolor{pgrey}]{$$},
  moredelim=[is][\textcolor{pgrey}]{\%\%}{\%\%},
  captionpos=b
}

\lstdefinestyle{basic}{  
  basicstyle=\footnotesize\ttfamily,
  breaklines=true
  numbers=left,
  numberstyle=\tiny\color{gray}\ttfamily,
  numbersep=7pt,
  backgroundcolor=\color{white},
  showspaces=false,
  showstringspaces=false,
  showtabs=false,
  frame=single,
  rulecolor=\color{black},
  captionpos=b,
  keywordstyle=\color{blue}\bf,
  commentstyle=\color{gray},
  stringstyle=\color{green},
  keywordstyle={[2]\color{red}\bf},
}


\lstdefinelanguage{custom}
{
morekeywords={public, void},
sensitive=false,
morecomment=[l]{//},
morecomment=[s]{/*}{*/},
morestring=[b]",
}


\lstdefinestyle{BashInputStyle}{
  language=bash,
  showstringspaces=false,
  basicstyle=\small\sffamily,
  numbers=left,
  numberstyle=\tiny,
  numbersep=5pt,
  frame=trlb,
  columns=fullflexible,
  backgroundcolor=\color{gray!20},
  linewidth=0.9\linewidth,
  xleftmargin=0.1\linewidth
}

%Logo in the upper right just change if you know what you are doing^^
\addtobeamertemplate{frametitle}{}{%
\begin{tikzpicture}[remember picture,overlay]
\node[anchor=north east,yshift=2pt] at (current page.north east) {\includegraphics[height=1.8cm]{htw}};
\end{tikzpicture}}

\begin{document}
\bibliographystyle{alpha}
\title{Netzwerke und verteilte Systeme\\
Übung Wintersemester 2022/21}
\subtitle{Organisatorisches\\

\href{mailto:Benjamin.Troester@HTW-Berlin.de}{Benjamin.Troester@HTW-Berlin.de}\\
		PGP: ADE1 3997 3D5D B25D 3F8F 0A51 A03A 3A24 978D D673 }

\author{Benjamin Tröster}

\date{}

\begin{frame}
\titlepage
\end{frame}

\section*{Road-Map}
\begin{frame}
\frametitle{Road-Map}
\begin{multicols}{2}
  \tableofcontents
\end{multicols}
\end{frame}

\section{Disclamer}
\begin{frame}
	\frametitle{Disclamer I}
	\begin{itemize}
		\item Objektive Kritik ist immer willkommen!
		\item Wenn sie Fehler oder Anmerkungen haben, tragen sie diese ruhig vor.
		\item Sie wollen ehrliches Feedback genauso wie ich
		\item Wenn Sie Anregungen, Verbesserungsvorschläge haben, nehme ich diese ebenfalls gerne an.
	\end{itemize}
\end{frame}

\section{Organisatorisches}
\subsection{Grundsätzliches}
\begin{frame}
	\frametitle{Grundsätzliches}
	\begin{itemize}
		\item Übungen: ab Oktober 2021 wahrscheinlich im Labor (wer weiß)
		\item (Mglw.: Meeting/Übung via Moodle: Telco-Plugin \emph{BigBlueButton (BBB)})
		\begin{itemize}
			\item Jede Übung hat eine eigene BBB-Session (unten bei den Übungsaufgaben)
		\end{itemize}
		\item Übungsgruppe \& -zeit $\Rightarrow$ LSF
		\begin{itemize}
			\item Beachten sie die Termine im LSF
		\end{itemize}
		\item Übungsblätter, Folien, Literatur \& Links etc. $\Rightarrow$ \url{moodle.htw-berlin.de}
		\item Zweiwöchentlicher Tonus 
		\begin{itemize}
			\item \textasciitilde 6 Übungsblätter -- bestehend aus Theorie- \& Laborteil
			\item Jede Übung in sich abgeschlossen
			\item \textbf{Aber:} Übungen bauen aufeinander auf
		\end{itemize}
	\end{itemize}
\end{frame}

\begin{frame}
	\frametitle{Klausur/Mündliche Prüfung}
	\begin{itemize}
		\item How to Pass:
		\begin{itemize}
			\item Dieses Semester: Es gibt aufgrund der hohen Durchfallquote wieder eine Klausurzulassung
			\begin{itemize}
				\item Themengebiete: Statisches Routing, NAT, DNS, DHCP, Wireshark
				\item Alle Themenkomplexe müssen lauffähig sein
				\item Änderung und kurze Rückfragen (Review der Aufgaben)
				\item Review-Termine: Werden in den Übungen bekanntgegeben
			\end{itemize}
			\item Es gibt jedoch keine Anwesenheitspflicht
			\item Ende des Semesters: mündliche Prüfung (50\% Übung, 50\% Vorlesung)
			\item Mündliche Prüfung: sehr praxisnah, klare Verbindung von VL und Ü
			\begin{itemize}
				\item Dauer 30 -45 Minuten, in Gruppen von 4-5 Studierenden
			\end{itemize}
			\item Das letzte Semester zeigt, dass sich die Anwesenheit in der Übung lohnt!
		\end{itemize}
	\end{itemize}
\end{frame}

\begin{frame}
	\frametitle{Übungsablauf}
	\begin{itemize}
		\item Beginn der Übung Nachbesprechung:
		\begin{itemize}
			\item Vorlesung
			\item Letzte Übung
		\end{itemize}
		\item Besprechung der Hausaufgaben
		\begin{itemize}
			\item Allgemeine Frage im Plenum für alle
			\item Fragen zum Laborübungsblatt
		\end{itemize}
		\item Umsetzung: Laboraufgaben
		\begin{itemize}
			\item Aufteilung in Kleingruppen via Breakout-Rooms
		\end{itemize}
	\end{itemize}
\end{frame}

\subsection{ECTS}
\begin{frame}
	\frametitle{Arbeitsaufwand -- ECTS}
	\begin{itemize}
		\item Modul Netzwerke: 5 ECTS (European Credit Transfer System) -- manchmal auch LP/CP (Leistungspunkte)
			\begin{itemize}
				\item 1 ECTS $\widehat{=}$ 30h
				\item Workload Netzwerke: 150h/Semester oder 37,5h/Monat oder $\sim$ 9,375h/Woche
				\item 2 SWS \footnote{Semesterwochenstunden} VL + 2 SWS Übung $\rightarrow$ 4 SWS oder 3h
				\item D.h. restlichen \textbf{6h/Woche}
				\begin{itemize}
					\item Vorbereitung \& Nachbereitung der Vorlesung
					\item Hausaufgaben, Recherche...
				\end{itemize}
			\end{itemize}
			\item Total Workload pro Semester: 30 CP $\widehat{=}$ 900h $\rightarrow$  $\sim$56.25/Woche
	\end{itemize}
\end{frame}
\end{document}